\scanpage{65}
% loc: 1–9, top
Soit $A$ une algèbre d'Azumaya de rang $4$ sur anneau $k$, et supposons que $\deleted{\unsure{A}}M$ soit ``triviale'', plus précisément, qu'elle provienne d'un $k$-module projectif de rang $2$ $E$, tel que le module inversible $\det E$ soit \unsure{trivial|triviale} (i.e. le \deleted{groupe structural} $\unsure{G_p(1)|\mathrm{GP}(1)|Gp(1)}_{\deleted{\unsure{k}}\,k}$-torseur décrivant $M$ peut se remonter en un $\mathrm{SL}(2)_k$-torseur.) Mais pour le moment, supposons seulem\textsuperscript{t} que $M$ soit triviale i.e.

% loc: 10
\[
 M\simeq\operatorname{End}(E).
\]

% loc: 11–14
Soit $u\in M$ un élément non dégénéré de carré nul. Alors $ku$ définit une ``droite'' $L\subset M$, et $D=ku\simeq L\otimes L'$, où $L'$ est l'orthogonal de $L$ dans $E^\vee\simeq E\otimes\deleted{\unsure{\det E}}\underline{\omega}^{-1}$ (où $\underline{\omega}=\omega(E)$).

% loc: 15–17
\deleted{Donc} En fait, sous l'identification précédente (dépendant d'une convention de signe) on aura $L'=L\otimes\underline{\omega}^{-1}$ donc

% loc: 18
\[
 D_L\simeq L^{\otimes2}\otimes\underline{\omega}^{-1},\quad
 \text{i.e.}\quad\underline{\omega}\xrightarrow[\sim]{\varphi_L}
 \underline{\operatorname{Hom}}(L^{\otimes2},D_L)
\]

% loc: 19–21
(isom can, \deleted{\unsure{à}} une fois faite la conv. de signe). Le choix d'une base $e\in\underline{\omega}$, d'où $e^{-1}\in\underline{\omega}^{-1}$ définit un iso. :

% loc: 22
\[
 \varphi_L^e:\quad\deleted{\unsure{D_L\longrightarrow{}}}\ L^{\otimes2}\xrightarrow{\sim} D_L
\]

% loc: 23–24
Si on remplace $e$ par $\lambda e$, $\unsure{e^{-1}|x^{-1}}$ est remplacé par $\lambda^{-1}e^{-1}$ et on trouve

% loc: 25
\[
 \varphi_L^{\lambda e}=\lambda\varphi_L^e
\]

% loc: 26, bottom
Quand $e$ est choisi (str. symplectique sur $E$)
% archive sheet number: 64

